\documentclass[10pt,a4paper]{amsart}
\usepackage[margin=19mm]{geometry}
\usepackage{amsmath,amssymb,mathtools}
\usepackage[scr=rsfs,cal=euler]{mathalfa}
\usepackage{xfrac}
\usepackage[unicode,hidelinks,bookmarks=false]{hyperref}
\newcommand{\ie}{i.e.\ }
\newcommand{\cf}{cf.\ }
\newcommand{\resp}{respectively\ }
\newcommand{\Subsection}{\subsection}
\providecommand{\nobreakdash}{\nobreak\hbox{-}}
\setlength{\emergencystretch}{2em}
\multlinegap0pt
\usepackage{amssymb}
\usepackage{mathtools}
\usepackage[shortlabels]{enumitem}
\usepackage{tikz}
\usepackage{xcolor}
\newtheorem{proposition}{Proposition}
\newcommand{\AlmostEnt}{\mathrm{AE}}
\title{\textbf{Folded-Algebraic Matroids: Characteristic Rigidity and Almost-Entropic Separation}}
\author{Shahram Khazaei}
\date{Source: arXiv:2609.01664v1 (CC BY 4.0)}
\begin{document}
\maketitle

\noindent\emph{Source and licence.} \textbf{Folded-Algebraic Matroids: Characteristic Rigidity and Almost-Entropic Separation}. Version arXiv:2609.01664v1 (CC BY 4.0), distributed by arXiv under the Creative Commons Attribution 4.0 International licence, http://creativecommons.org/licenses/by/4.0/, as recorded in the arXiv licence metadata for this exact version and shown on the abstract page https://arxiv.org/abs/2609.01664v1. Full licence notice: \texttt{licenses/arxiv-2609.01664-CC-BY-4.0.txt}.\par\smallskip

\noindent\textit{Editorial scope.} Counted unit: the article's proposition prop:welding (source0.tex lines 989--997). The article's complete original proof of it is delivered with it (source0.tex lines 998--1026, one \texttt{proof} environment of 29 lines). No mathematics is written, repaired or re-derived: the statement and the proof are the article's own lines, in the article's own notation, and no retained line is substituted or re-flowed.  No further source-local context is needed: every cross-reference of every retained line resolves inside the two delivered spans.  Nothing was omitted from any retained span, so the package carries no omission record.  No retained line contains a citation, so the wrapper carries no bibliography and no bibliography entry.  No retained line contains a figure environment, an image-inclusion command or drawing code, so no image is delivered and the package declares no asset dependency.  Editorial formatting only, in addition: an AMS \texttt{amsart} preamble that restates the article's own declarations and macros used by the retained lines, namely the environments \texttt{proposition} on separate counters, together with the standard packages those lines need.  The retained environments print in this document as Proposition 1 in order of appearance, whereas the article prints the same items with the numbering of its own full document.  The complete native source of the article as distributed in the arXiv e-print for this exact version is bundled unchanged as \texttt{sources/209132/source0.tex}.
\begin{proposition}[Independent-set welding]\label{prop:welding}
Let $M_1,M_2\in\AlmostEnt$ have disjoint ground sets, and choose ordered independent sets
\[
I_1=(x_1,\dots,x_k),\qquad I_2=(y_1,\dots,y_k).
\]
Starting from $M_1\oplus M_2$, identify $x_j$ with $y_j$ successively by the preceding principal-extension--contraction--deletion operation.  The resulting matroid $W$ is almost entropic and contains natural restrictions isomorphic to both $M_1$ and $M_2$.

If $I_1$ and $I_2$ are bases of two rank-$r$ matroids and $k=r$, then $r(W)=r$ and both factor restrictions are spanning.
\end{proposition}

\begin{proof}
Almost-entropicity follows from closure under direct sums, principal extensions, contractions, and deletions.  It remains to check that neither factor is damaged by the identifications.

Keep the original disjoint labels while performing the construction and, after the last deletion, retain $x_j$ as the common label.  For $S$ in the final ground set an induction on the $k$ identifications gives
\begin{equation}\label{eq:weld-rank}
 r_W(S)=
 \min_{J\subseteq[k]}
 \left(
 r_{M_1\oplus M_2}
 \bigl(S\cup\{x_j,y_j:j\in J\}\bigr)-|J|
 \right).
\end{equation}
Indeed, one identification transforms a rank function by
\[
 r(S)\longmapsto \min\{r(S),r(Sxy)-1\},
\]
and the resulting operators commute for the disjoint pairs.

If $S\subseteq E(M_1)$, then the $M_2$-contribution of the adjoined $y_j$ is $|J|$ because $I_2$ is independent.  Every term of \eqref{eq:weld-rank} is therefore at least $r_{M_1}(S)$, while $J=\varnothing$ gives equality.  Thus $M_1$ survives as a restriction.

Now let $T\subseteq E(M_2)$ and let $S$ be its image after replacing every selected $y_j$ by the common label $x_j$.  Put $I=\{j:y_j\in T\}$.  Taking $J=I$ in \eqref{eq:weld-rank} gives $r_{M_2}(T)$.  For arbitrary $J$, independence of $I_1$ contributes $|I\cup J|$, while deleting the elements $y_j$ with $j\in I\setminus J$ can lower $r_{M_2}(T)$ by at most $|I\setminus J|$.  Hence the corresponding term is at least
\[
 |I\cup J|+r_{M_2}(T)-|I\setminus J|-|J|
 =r_{M_2}(T).
\]
So $M_2$ also survives exactly.

If both chosen sets are full bases of rank $r$, the surviving $M_1$ restriction gives $r(W)\ge r$, while taking $J=[r]$ in \eqref{eq:weld-rank} for the whole ground set gives $r(W)\le r+r-r=r$.
\end{proof}

\end{document}
